This section describes the traceability between the system requirements and the software components responsible for their implementation. Each requirement defined in Section~3 of the RASD document is mapped to the manager that fulfills it, we also added a Client Application section to highlight sensor-based requirements.

\section{Functional Requirements}
\subsection*{Authentication Manager}
\begin{itemize}
	\item \textbf{R1}: The system shall allow new users to register an account
	\item \textbf{R2}: The system shall allow users not registered to access the application without an account
	\item \textbf{R3}: The system shall allow registered users to log in to the application
	\item \textbf{R5}: The system shall allow users to log out of their current session
\end{itemize}

\subsection*{Profile Manager}
\begin{itemize}
	\item \textbf{R4}: The system shall allow registered users to view and update their personal profile information
	\item \textbf{R47}: The system shall distinguish between declared public trips and private ones
	\item \textbf{R48}: The system shall allow registered users to visualize a summary of the trip
\end{itemize}

\subsection*{Recording Manager}
\begin{itemize}
	\item \textbf{R6}: The system shall allow registered users to record their cycling trips
	\item \textbf{R7}: The system shall store recorded trips to maintain a history of the user’s cycling activities
	\item \textbf{R8}: The system shall compute and display total distance covered for each trip
	\item \textbf{R9}: The system shall compute and display average speed for each trip
	\item \textbf{R10}: The system shall compute and display total time spent for each trip
	\item \textbf{R11}: The system shall compute and display elevation gain for each trip
	\item \textbf{R12}: The system shall compute and display calories spent for each trip
	\item \textbf{R13}: The system shall compute and display additional performance metrics for the trips
	\item \textbf{R28}: The system shall analyze sensor data to identify potential path irregularities
	\item \textbf{R46}: The system shall allow registered users to manually insert the duration of the trip
\end{itemize}

\subsection*{Storage Manager}
\begin{itemize}
	\item \textbf{R7}: The system shall store recorded trips to maintain a history of the user’s cycling activities
	\item \textbf{R38}: The system shall collect paths data and keep it saved
\end{itemize}

\subsection*{Look Publish Manager}
\begin{itemize}
	\item \textbf{R15}: The system shall allow registered users to insert information about bike paths
	\item \textbf{R16}: The system shall allow registered users to make publishable information about bike paths
	\item \textbf{R17}: The system shall allow registered users to add info about the status of a bike path
	\item \textbf{R18}: The system shall allow registered users to add the presence of relevant obstacles on the bike path
	\item \textbf{R31}: The system shall allow registered users to mark their inserted data as ``publishable''
	\item \textbf{R39}: The system shall visualize paths data
	\item \textbf{R40}: The system shall filter paths through various attributes
	\item \textbf{R41}: The system shall allow the user to search through public paths
\end{itemize}

\subsection*{Search Manager}
\begin{itemize}
	\item \textbf{R21}: The system shall calculate the distance of a path
	\item \textbf{R34}: The system shall find bike paths given an origin and a destination
	\item \textbf{R35}: The system shall visualize the bike paths connecting the specified origin and destination on a map
	\item \textbf{R36}: The system shall compute a score for each path based on its status and effectiveness
	\item \textbf{R37}: If multiple paths exist, the system shall display them sorted by their calculated score
\end{itemize}

\subsection*{Merge Manager}
\begin{itemize}
	\item \textbf{R32}: The system shall merge information about the same path from multiple users
	\item \textbf{R33}: The system shall prioritize data based on timeframe and number of confirmations
\end{itemize}

\subsection*{Notification Manager}
\begin{itemize}
	\item \textbf{R30}: The system shall require the user to confirm or correct automatically acquired data
\end{itemize}

\subsection*{Client Application}
\begin{itemize}
	\item \textbf{R6}: The system shall allow registered users to record their cycling trips
	\item \textbf{R25}: The system shall automatically infer that the registered user is cycling based on their speed
	\item \textbf{R26}: The system shall collect GPS data to reconstruct the path followed by the user during the ride
	\item \textbf{R27}: The system shall acquire data from the mobile device’s accelerometer and gyroscope
	\item \textbf{R29}: The system shall require the user to confirm or correct automatically acquired sensor data
\end{itemize}

\section{Non-Functional Requirements}
As previously described in the RASD document in section~3.3, here we explain how the performance requirements are granted through design decisions.

\subsection{Latency and Concurrency}
By having standard 4G/5G connection and a solid client-side application, it is sufficient to send numerical coordinates as data from the server to the application to locally render the path on the map visualization of the device, also as stated in the RASD document the publication of path doesn't need much work of the client-side device and can be done asynchronously. The large number of users is granted by the possibility of having virtually as many application servers as the owner of the system wants thanks to the implemented architecture and the API Gateway.

\subsection{Data Storage}
Having a main database as the head of write operations coupled with read-only databases used for historical data (such as old trips and paths) and a local client-side cache made for personal trips only and immediatly accessible data (such as trip statistics) meets the standard described in the RASD.